\documentclass{article}

\PassOptionsToPackage{numbers, compress}{natbib}
\usepackage[main, final]{neurips_2026}

\usepackage[utf8]{inputenc}
\usepackage[T1]{fontenc}
\usepackage{hyperref}
\usepackage{url}
\usepackage{booktabs}
\usepackage{amsfonts}
\usepackage{nicefrac}
\usepackage{microtype}
\usepackage{xcolor}

% remove the conference footer line on page 1
\makeatletter
\renewcommand{\@noticestring}{}
\makeatother

\title{Evidence-Aware Financial News Forecasting}

\author{%
  Vyom Shah \\
  Western University \\
  \texttt{vshah259@uwo.ca} \\
  \And
  Harshil Patel \\
  Western University \\
  \texttt{hpate566@uwo.ca} \\
  \And
  Harmit Patel \\
  Western University \\
  \texttt{hpate578@uwo.ca} \\
}

\begin{document}

\maketitle

\begin{abstract}
Financial news can change investors' expectations, yet a sentiment score does not show whether an article is new, whether the market has already reacted, or whether news improves a forecast over prices. We propose SignalProof, a reproducible research framework and evidence-aware financial intelligence product. Using FNSPID and market prices, we will measure what news attention, sentiment, and event information add to next-day large-price-move forecasts. A second-stage model will estimate when news is likely to help, and can keep the price-only forecast or withhold an unsupported signal. Chronological testing, calibrated probabilities, and an auditable forecast record will link the research to an analyst dashboard. We will publish the protocol and results, including null findings.
\end{abstract}

\section{Introduction}
Analysts see many articles about one company, from repeated coverage to new information, and tone alone cannot separate them. Prior studies connect media tone to market behavior \citep{tetlock2007}, but apparent gains can come from poor timestamp alignment, repeated articles, or retrospective model selection. Our aim is not profitable trading. It is to determine when, and with what uncertainty, news adds information about tomorrow's stock-movement risk. We test three questions: (H1) do sentiment and event features improve large-move probability forecasts beyond prices and article volume; (H2) can a conditional model identify cases in which news reduces out-of-sample forecast loss; and (H3) do those benefits survive time, sector, coverage, and market-regime shifts? We test whether news has measurable incremental predictive value beyond a matched price-only model.

\section{Proposal}
\paragraph{Data and auditable information timing.} FNSPID \citep{dong2024} links historical financial news and stock information. We will audit timestamps, ticker links, duplication, missingness, and coverage before selecting a stratified company sample. Adjusted daily OHLCV with market and sector benchmarks provide price-only controls. At each prediction cutoff (US/Eastern close), features use only information demonstrably available then. We will handle holidays, uncertain timestamps, and survivorship.

\paragraph{From articles to evidence.} We will compare FinBERT \citep{araci2019}, VADER \citep{hutto2014}, and the financial lexicon of Loughran and McDonald \citep{loughran2011}. We will aggregate polarity, disagreement, article count, and trailing changes by ticker and date, and cluster related headlines into events, validating labels and novelty on a manually reviewed sample. Price features include lagged returns, momentum, volatility, volume, and market conditions. The primary label is a next-session absolute return above a 90th-percentile threshold learned from training data; secondary analyses cover direction and longer horizons.

\paragraph{Models, hypotheses, and evaluation.} We will compare a historical-rate baseline, price-only, price+attention, price+sentiment, and price+attention+event/sentiment models using regularized regression and LightGBM \citep{ke2017}. For H2, an incremental-value gate will learn from earlier out-of-fold losses whether to use the price-only model, the news-enhanced model, or abstain, with simple gating rules as controls at matched coverage. Selective forecasting is not new; our contribution is empirical. Nested, purged expanding-window tests with past-only tuning will assess Brier score (primary), AUROC, calibration, and risk-coverage curves. Paired date-block bootstrap intervals and corrected secondary tests will address dependence and multiple comparisons. A final untouched period checks generalization.

\paragraph{Chance, association, and causation.} Every gain carries a bootstrap confidence interval, so readers can judge whether it could be chance, and corrected tests limit false discoveries. A reliable gain would show only association, not causation: an article may report a move already under way, and both may respond to one event. We make no causal claims and reduce this risk with strict cutoffs, price-only controls, and validated event labels.

\paragraph{Visualization plan.} We will plot (1) Brier-score differences versus price-only with confidence intervals, (2) calibration curves, (3) risk-coverage curves, (4) large-move rate by sentiment decile, and (5) event timeline.

\paragraph{Research output and product.} For analysts who must tell important developments from recycled coverage, SignalProof will show a stock/event timeline, de-duplicated source-linked coverage, price-only versus news-enhanced risk estimates, model calibration, evidence sufficiency, and a log of when forecasts were issued and how they performed. A localhost prototype may later become a deployed application with scheduled ingestion, an API, watchlists, alerts, and monitoring, subject to live-feed rights. It supports research and risk awareness, not individualized investment advice. Code and protocol will be released where permitted.

\section{Related Work}
Tetlock \citep{tetlock2007} studies media pessimism and prices; Loughran and McDonald \citep{loughran2011} motivate finance-specific language analysis; FinBERT \citep{araci2019} brings pretrained financial NLP. FNSPID \citep{dong2024} already reports modest gains from adding sentiment to forecasting; Lopez-Lira and Tang \citep{lopezlira2023} investigate language-model headline signals. Selective, uncertainty-aware financial prediction exists, including FinAbstain \citep{torres2026}, so abstention alone is not our novelty claim. We instead measure conditional incremental value over a matched price-only baseline, with event/attention ablations and audited information cutoffs. Our experiments will test whether this improves on alternatives.

\begin{thebibliography}{99}
\bibitem{araci2019} Dogu Araci. FinBERT: Financial sentiment analysis with pre-trained language models. \emph{arXiv:1908.10063}, 2019.
\bibitem{dong2024} Zihan Dong, Xinyu Fan, and Zhiyuan Peng. FNSPID: A comprehensive financial news dataset in time series. \emph{arXiv:2402.06698}, 2024.
\bibitem{hutto2014} Clayton J. Hutto and Eric Gilbert. VADER: A parsimonious rule-based model for sentiment analysis of social media text. \emph{Proc. AAAI Int. Conf. on Web and Social Media}, 2014.
\bibitem{ke2017} Guolin Ke, Qi Meng, Thomas Finley, Taifeng Wang, Wei Chen, Weidong Ma, Qiwei Ye, and Tie-Yan Liu. LightGBM: A highly efficient gradient boosting decision tree. \emph{Advances in Neural Information Processing Systems}, 30, 2017.
\bibitem{lopezlira2023} Alejandro Lopez-Lira and Yuehua Tang. Can ChatGPT forecast stock price movements? \emph{arXiv:2304.07619}, 2023.
\bibitem{loughran2011} Tim Loughran and Bill McDonald. When is a liability not a liability? \emph{The Journal of Finance}, 66(1), 2011.
\bibitem{tetlock2007} Paul C. Tetlock. Giving content to investor sentiment: The role of media in the stock market. \emph{The Journal of Finance}, 62(3), 2007.
\bibitem{torres2026} Dorothy Torres, Wei Cheng, and Henan Huang. FinAbstain: Uncertainty-calibrated multimodal RAG for selective financial forecasting. \emph{arXiv:2607.24875}, 2026.
\end{thebibliography}

\end{document}
